\begin{theorem}
Let \(s,n\ge 3\), let \(k\ge 1\), and suppose \(F\) and \(G\) are \((n,d(F),\lambda(F))\)- and \((n,d(G),\lambda(G))\)-graphs on the same vertex set, with \(d(G)>0\). Assume \((F,G)\) is \(H_s\)-free, and set
\[
  \eta=\max\left\{\frac{\lambda(G)^2}{d(G)^2},
  \frac{\lambda(F)\lambda(G)}{d(F)d(G)}\right\},
  \qquad
  w=\frac{4n\log n}{d(G)}.
\]
Assume also that \(n^{-2}\le \eta\le 1\).
If \(w\le k\le \eta d(F)n\), then
\[
  r(s,k)\ge \frac{1}{50} k\,\eta^{(w-k)/k}-1.
\]
In particular, if \(\eta d(F)n\ge k\ge 100n(\log n)^2/d(G)\), then
\[
  r(s,k)\ge \frac{k}{100\eta}-1.
\]
\end{theorem}
\begin{proof}
Let
\[
  M=d(F)n.
\]
Assume first that \(w\le k\le \eta M\). By \(\noderef{ProductDigraphForwardIndependentBound}\), using the hypotheses \(s,n\ge3\), there is a \(T_s\)-free digraph \(D\) on \(M\) vertices such that
\[
  \overrightarrow{i_k}(D)\le 16^k\eta^{k-w}M^k.
\]
Since \(k\ge1\), \(\noderef{DigraphToGraphIndependentSetBound}\) applies to \(D\). It gives a \(K_s\)-free graph \(\Gamma\) on the same \(M\)-vertex set and satisfying
\[
  i_k(\Gamma)\le \left(\frac{e}{k}\right)^k\overrightarrow{i_k}(D).
\]
Combining the two displayed inequalities gives
\[
  i_k(\Gamma)\le
  \left(\frac{e}{k}\right)^k16^k\eta^{k-w}M^k. \tag{1}
\]

Define
\[
  B=\left(\frac{e}{k}\right)^k16^k\eta^{k-w}M^k,
  \qquad
  p=B^{-1/k}.
\]
Because \(n\ge3\), the lower bound \(n^{-2}\le\eta\) implies \(\eta>0\). The antecedent \(k\le \eta M\), together with \(k\ge1\), therefore implies \(M>0\), so \(B>0\) and hence \(p>0\). From (1) we have \(i_k(\Gamma)\le B\), and by the definition of \(p\),
\[
  p^k B=1.
\]
Therefore
\[
  p^k i_k(\Gamma)\le 1. \tag{2}
\]

We next check \(p\le1\). Expanding the definition of \(p\) gives
\[
  p=\frac{k}{16eM}\,\eta^{(w-k)/k}.
\]
Because \(d(G)>0\), \(n\ge3\), and \(w=4n\log n/d(G)\), we have \(w>0\). Together with \(w\le k\), this gives
\[
  -1\le \frac{w-k}{k}\le0.
\]
Since \(0<\eta\le1\), it follows that
\[
  \eta^{(w-k)/k}\le \eta^{-1}.
\]
Using \(k\le \eta M\), we obtain
\[
  p\le \frac{k}{16eM}\eta^{-1}\le \frac1{16e}<1.
\]
Thus \(0\le p\le1\).

Now \(\Gamma\) is \(K_s\)-free, \(k\ge1\), \(0\le p\le1\), and (2) is the hypothesis \(p^k i_k(\Gamma)\le1\). Applying \(\noderef{SamplingKsFreeRamseyBound}\) gives
\[
  r(s,k)>pM-1.
\]
Using the formula for \(p\), this becomes
\[
  r(s,k)>
  \frac{k}{16e}\eta^{(w-k)/k}-1.
\]
Since \(16e<50\), the right hand side is at least
\[
  \frac{k}{50}\eta^{(w-k)/k}-1.
\]
Therefore
\[
  r(s,k)\ge \frac{1}{50}k\eta^{(w-k)/k}-1,
\]
which proves the first conclusion.

For the ``in particular'' statement, assume
\[
  k\le \eta M
  \qquad\text{and}\qquad
  k\ge \frac{100n(\log n)^2}{d(G)}.
\]
The first conclusion applies once we know \(w\le k\). This follows from the second inequality for all \(n\ge3\): since \(\log n>0\) and \(w>0\),
\[
  \frac{k}{w}
  \ge
  \frac{100n(\log n)^2/d(G)}{4n\log n/d(G)}
  =25\log n>1.
\]
Thus \(w\le k\). Applying the first conclusion gives
\[
  r(s,k)\ge \frac{k}{50}\eta^{(w-k)/k}-1
  =\frac{k}{50\eta}\eta^{w/k}-1. \tag{3}
\]

It remains to bound \(\eta^{w/k}\) from below. From \(n^{-2}\le\eta\le1\) and \(w/k\ge0\),
\[
  \eta^{w/k}\ge (n^{-2})^{w/k}.
\]
The lower bound on \(k\), together with the definition of \(w\), gives
\[
  \frac{w}{k}\le \frac{1}{25\log n}.
\]
Because \(n^{-2}\in(0,1)\), decreasing the exponent increases the value, so
\[
  (n^{-2})^{w/k}\ge (n^{-2})^{1/(25\log n)}
  =\exp(-2/25).
\]
Since \(\exp(-2/25)>\exp(-1/10)>1/2\), we have
\[
  \eta^{w/k}>1/2.
\]
Substituting this into (3) yields
\[
  r(s,k)>\frac{k}{100\eta}-1.
\]
In particular,
\[
  r(s,k)\ge \frac{k}{100\eta}-1,
\]
which is the claimed second bound.
\end{proof}
